%! AP Statistics — Unit 3, Lesson 3.1 — Guided Notes KEY
\documentclass[10pt]{article}
\usepackage{apstats-article}
\usepackage{apstats-key}

\begin{document}

\pageheader{Unit 3, Lesson 3.1}{Guided Notes — KEY}
\namedateperiod

\vspace{0.12in}

\begin{objectivebox}
\small By the end of this lesson, I will be able to\dots\\[4pt]
\ans{Explain why non-chance methods produce untrustworthy conclusions.}\\[1pt]
\ans{Explain why a large non-random sample is not more trustworthy than a small random one.}\\[1pt]
\ans{Identify whose opinions are — and are not — represented in a data collection scenario.}
\end{objectivebox}

\vspace{0.1in}

\begin{vocabbox}
\termblanklong{Data collection method} \ans{The procedure used to gather data from a population or sample.}
\termblanklong{Chance-based method} \ans{A method that uses randomness to select individuals, giving each a known probability of selection.}
\termblanklong{Non-chance method} \ans{A method that does not use randomness; convenience, voluntary response, or judgment.}
\termblanklong{Bias} \ans{A systematic tendency for the sample to over- or under-represent certain responses.}
\termblanklong{Anecdotal evidence} \ans{Data based on personal stories or isolated examples rather than systematic collection.}
\end{vocabbox}

\vspace{0.1in}

\begin{hookbox}
\small\textbf{Three surveys:}
\begin{enumerate}[leftmargin=1.5em, itemsep=4pt]
  \item A grocery store comment card near the exit: 12 responses, all complaints.
  \item A news website polls online visitors about a school bond.
  \item A researcher randomly dials 400 phone numbers.
\end{enumerate}
\textbf{Which would you trust most? Why?}\\[4pt]
\ans{Survey 3 — random dialing gives every household a chance to be selected; the other two rely on voluntary response (only people with strong opinions respond), which produces a biased sample.}
\end{hookbox}

\vspace{0.1in}

\begin{notesbox}{The Trustworthiness of Data}
\small
\textbf{Key Principle (VAR-1.E.1):} Methods for data collection that do \textbf{not}
rely on \ans{chance} result in \ans{untrustworthy} conclusions.

\vspace{10pt}
\renewcommand{\arraystretch}{1.6}
\begin{tabularx}{\linewidth}{@{} >{\bfseries\color{navy}}p{0.3\linewidth} X @{}}
Chance-based method & \ans{Randomness is used; every individual has a known probability of being selected.} \\
Non-chance method & \ans{No randomness; individuals are chosen by convenience, self-selection, or judgment.} \\
\end{tabularx}

\vspace{10pt}
\textbf{Counter-intuitive fact:} A \ans{large} non-random sample is \textbf{not} more
trustworthy than a \ans{small} random sample.

\vspace{8pt}
\textbf{The 1936 Literary Digest Poll:}
\begin{itemize}[itemsep=2pt]
  \item Contacted \ans{10 million} people --- one of the largest polls ever.
  \item Predicted \ans{Alf Landon} would win the presidential election.
  \item Actual winner: \ans{Franklin D. Roosevelt}
  \item Why it failed: \ans{The sample came from car registrations and phone directories, excluding poor voters who overwhelmingly supported Roosevelt.}
\end{itemize}
\end{notesbox}

\vspace{0.1in}

\begin{notesbox}{Evaluating a Data Collection Method}
\small

\vspace{6pt}
\renewcommand{\arraystretch}{1.8}
\begin{tabularx}{\linewidth}{@{} l X @{}}
\textbf{Question} & \textbf{My answer} \\
\hline
Was chance used to select individuals? & \ans{Look for words like "randomly selected," "SRS," or "random sample."} \\
Whose opinions/data are represented? & \ans{The people who responded — which may not be the target population.} \\
Who is \emph{not} represented? & \ans{Anyone excluded by the sampling frame or who chose not to respond.} \\
Can conclusions be trusted for the whole population? & \ans{Only if chance was used; otherwise, conclusions are limited to those who responded.} \\
\end{tabularx}
\end{notesbox}

\vspace{0.1in}

\begin{practicebox}
\small
\vspace{6pt}
\begin{enumerate}[leftmargin=1.5em, itemsep=10pt]
  \item A radio station invites listeners to text in their opinion about a new law.
        \ans{$\times$} \quad \ans{Voluntary response — only those with strong opinions will text.}

  \item A researcher uses a random number table to select 200 voters from the registered voter list.
        \ans{\checkmark} \quad \ans{Chance-based — every registered voter has an equal chance of being selected.}

  \item A teacher surveys her own class to find out how all students in the school feel about homework.
        \ans{$\times$} \quad \ans{Convenience sample — only her students are represented; they may not reflect the whole school.}
\end{enumerate}
\end{practicebox}

\end{document}
